\documentclass{article}
%
\newcommand{\cat}[1]{\mathsf{#1}}
\newcommand{\univ}[1]{\mathcal{#1}}
\newcommand{\path}{\equiv}
\newcommand{\id}[1]{\mathsf{#1}}
\usepackage{amsmath}
\usepackage{amsthm}
\usepackage{amssymb}
\usepackage{tikz}
\usepackage{quiver}
\newtheorem{theorem}{Theorem}
\newtheorem{prop}[theorem]{Proposi{t}ion}
\newtheorem{lemma}[theorem]{Lemma}
\theoremstyle{plain}
\newtheorem{corollary}[theorem]{Corollary}
\theoremstyle{remark}
\newtheorem{defn}[theorem]{Def{i}ni{t}ion}
\usepackage{xcolor}
\usepackage{natbib}

% From CAN package "orcidlink"
\usepackage{scalerel}
\usepackage{tikz}
\usetikzlibrary{svg.path}
\definecolor{orcidlogocol}{HTML}{A6CE39}
\tikzset{
  orcidlogo/.pic={
    \fill[orcidlogocol] svg{M256,128c0,70.7-57.3,128-128,128C57.3,256,0,198.7,0,128C0,57.3,57.3,0,128,0C198.7,0,256,57.3,256,128z};
    \fill[white] svg{M86.3,186.2H70.9V79.1h15.4v48.4V186.2z}
                 svg{M108.9,79.1h41.6c39.6,0,57,28.3,57,53.6c0,27.5-21.5,53.6-56.8,53.6h-41.8V79.1z M124.3,172.4h24.5c34.9,0,42.9-26.5,42.9-39.7c0-21.5-13.7-39.7-43.7-39.7h-23.7V172.4z}
                 svg{M88.7,56.8c0,5.5-4.5,10.1-10.1,10.1c-5.6,0-10.1-4.6-10.1-10.1c0-5.6,4.5-10.1,10.1-10.1C84.2,46.7,88.7,51.3,88.7,56.8z};
  }
}

\newcommand\orcidicon[1]{\href{https://orcid.org/#1}{\mbox{\scalerel*{
\begin{tikzpicture}[yscale=-1,transform shape]
\pic{orcidlogo};
\end{tikzpicture}
}{|}}}}
% end orcidlink
\usepackage{hyperref}
% notation 
% end notationa
\begin{document}
\title{Regular and Exact Categories in Realisability}
\date{\today}
\author{Rahul Chhabra \orcidicon{0009-0007-7917-6461}\\
{Indian Institute of Information Technology\\}}
\maketitle
\begin{abstract}
    These notes are an exposition of internal (equivalence) relations, kernel pairs, regular epis, regular categories, effectiveness of internal relations, 
    exact categories, exact functors and exact completion of a regular category. All of these ideas are essential to the construction of realisability topoi. The category of assemblies over a 
    combinatory algebra $A$, $\cat{Asm}(A)$ is a regular category and the exact completion of $\cat{Asm}(A)$ is exactly the realisability topos over $A$, $\cat{RT}(A)$.
\end{abstract}

\section{Introduction}

TODO : Write proper introduction.

\section{Regular Categories}

\subsection{Kernel Pairs and the Image of a Function}

\subsubsection{Intution and Motivation}

The \textbf{image} of a map $f : A \to B$ is defined as 

\begin{equation*}
    \id{Im}(f) :\equiv \mathop{\Sigma}\limits_{y : B}^{} \mathop{\exists} \limits_{x : A}^{} f(x) = y 
\end{equation*}. 

In words, this is the total space of all $y : B$ such that there merely exists an $x : A$ for which we can witness $f x \path y$. 
Notice that the usage of mere existence means that even if $A$ and $B$ are arbitrary higher groupoids - whether or not a $y : B$ is in the image of $f$ is a \emph{property} and \emph{not structure}.

Notice how the definition suggests that the image should be a \emph{subobject} of the codomain. 
In the HoTT context, we must instead settle for saying that the image \emph{embeds into} the codomain. For a proof, see \citep{The_Agda_Community_Cubical_Agda_Library_2023}. 
The notion of embedding behaves the same as injections on the level of sets but is better behaved at higher levels. 

The most trivial inhabitants of the image are obtained if you already have an $x : A$. 

\begin{equation*}
    \id{inc}(x) :\equiv \; \mid (f (x) , x , \id{refl}(x)) \mid 
\end{equation*}

Let us take an illustrative example. Let $f : A \to B$ and $x_1 \; x_2 : A$ such that $p : f \; x_1 = f \; x_2$. 
We shall now prove that $\id{inc}(x_1) = \id{inc}(x_2)$.

\begin{lemma}
    \begin{equation*}
        \prod_{f : A \to B}^{} \prod_{x_1 \; x_2 : A}^{} \; \prod_{p : f x_1 \equiv f x_2} \; \id{inc}(x_1) = \id{inc}(x_2)    
    \end{equation*}
    
\end{lemma}

\begin{proof}
    We intend to construct a path between two inhabitants of a $\sum$ type. 
    Recall that two dependent pairs $(a, p_a)$ and $(b,p_b)$ have a path between them if there is a path between the independent components $j : a \equiv b$ and a path between the dependent components \emph{over} $j$. 

    In our case, the dependent component is a property, so it suffices to have a path between the independent components, which is given by $p$.
\end{proof}

Thus we present an important insight - inside the image of a function, it is not possible to distinguish between points that have the same output on the function.

\subsubsection{Kernel Pairs and Coequalisers}

Let us start by asking - what categorical structure could capture the intuition that it is not possible to distinguish between two kinds of points. 

A good candidate would be coequalisers. 

Recall that for a pair of arrows $f \; g : A \to B$ their coequaliser is the arrow $\id{coeq} : B \to C$ to an object $C$ such that $\id{coeq} \circ f = \id{coeq} \circ g$ and $C$ and $\id{coeq}$ are \emph{universal} so that if $q : B \to D$ is another arrow such that $q \circ f = q \circ g$ then $q$ factorises through $\id{coeq}$ uniquely. 


% https://q.uiver.app/#q=WzAsNCxbMCwwLCJBIl0sWzIsMCwiQiJdLFszLDAsIkMiXSxbMywxLCJEIl0sWzAsMSwiZiIsMCx7Im9mZnNldCI6LTJ9XSxbMCwxLCJnIiwyLHsib2Zmc2V0IjoyfV0sWzEsMiwiY29lcSJdLFsxLDMsInEiXSxbMiwzLCJcXGV4aXN0cyEiLDEseyJzdHlsZSI6eyJib2R5Ijp7Im5hbWUiOiJkYXNoZWQifX19XV0=
\[\begin{tikzcd}
	A && B & C \\
	&&& D
	\arrow["f", shift left=2, from=1-1, to=1-3]
	\arrow["g"', shift right=2, from=1-1, to=1-3]
	\arrow["coeq", from=1-3, to=1-4]
	\arrow["q", from=1-3, to=2-4]
	\arrow["{\exists!}"{description}, dashed, from=1-4, to=2-4]
\end{tikzcd}\]

Notice that "inside" the coequaliser $f$ and $g$ are indistinguishable. We wish to achieve something similar, except we would have pairs $(x_1,x_2)$ with $f(x_1) = f(x_2)$.

Firstly, we need to sort out the pair situation. This can be done using pullbacks. 

Consider the pullback of $f$ along $f$, that is, consider the following pullback square :

% https://q.uiver.app/#q=WzAsNCxbMSwwLCJBIl0sWzAsMSwiQSJdLFsxLDEsIkIiXSxbMCwwLCJBXFx0aW1lc19CQSJdLFsxLDIsImYiXSxbMCwyLCJmIiwyXSxbMywwLCJcXHBpXzIiLDJdLFszLDEsIlxccGlfMSJdXQ==
\[\begin{tikzcd}
	{A\times_BA} & A \\
	A & B
	\arrow["f", from=2-1, to=2-2]
	\arrow["f"', from=1-2, to=2-2]
	\arrow["{\pi_2}"', from=1-1, to=1-2]
	\arrow["{\pi_1}", from=1-1, to=2-1]
\end{tikzcd}\].

We call $\pi_1$ and $\pi_2$ the \textbf{kernel pairs} of $f$. 

What would the coequaliser of the kernel pairs look like? 
We would have an object such that "inside" the object one cannot distinguish between $x_1$ and $x_2$ if $f(x_1) = f(x_2)$. 

This is exactly we set out to achieve!

\subsection{Regular Categories}

We are now in a position to define regular categories. We simply abstract out all of these observations.

\begin{defn}
    A \textbf{regular category} \citep{borceux_1994} is a category with :
    \begin{enumerate}
        \item All finite limits 
        \item Kernel pairs of all morphisms
        \item Coequalisers of all kernel pairs 
        \item Stability of regular epics under pullback
    \end{enumerate}
\end{defn}

I have not been able to intuit why the last condition (stability of regular epics under pullback) is necessary, but it is likely either a coherence condition (to exclude pathological models) or has to do with the many things we would like to do with regular categories, such as interpret regular logic and construct the allegory of relations.

\bibliographystyle{plainnat}
\bibliography{realisability}

\end{document}
